% !TeX root = ../../relatorio-projeto-aplicado.tex
% Substitua integralmente as instrucoes por uma sintese do trabalho concluido.
\undfsetcontent{resumo}{%
Este texto orienta a redação do resumo de um relatório final de Projeto Aplicado.
Apresente o tema, a necessidade investigada, o contexto e o objetivo principal.
Delimite o projeto realizado no semestre, evitando confundi-lo com todos os
projetos que poderão integrar um futuro portfólio de trabalho de conclusão de
curso. Descreva os procedimentos efetivamente executados, as fontes ou materiais
utilizados e os critérios empregados para avaliar as entregas. Sintetize a
solução produzida, os resultados obtidos e a conclusão sustentada pelas
evidências disponíveis. Identifique o estágio real de desenvolvimento e
distinga testes técnicos de avaliação com pessoas ou de efeitos observados
no ambiente atendido. Se houve alteração de escopo, registre sua consequência
quando necessária à compreensão do resultado. Não apresente benefícios
esperados como resultados comprovados nem transforme o resumo em um plano
de atividades futuras. Redija um parágrafo único, com linguagem objetiva e
sem referências bibliográficas ou informações ausentes do relatório.
Substitua estas instruções pela síntese fiel do seu trabalho antes da entrega.
Confira as palavras-chave e a correspondência entre o resumo em português
e a versão em língua estrangeira, seguindo as orientações do curso.}
\undfsetcontent{palavraschave}{projeto aplicado; relatório final; avaliação de resultados}
\undfsetcontent{abstract}{%
This paragraph explains how to write the abstract of a completed applied
project report. Present the topic, identified need, context and main objective.
Define the scope of the project completed during the semester, keeping it
distinct from other projects that may later be selected for an undergraduate
portfolio. Describe the procedures actually performed, the sources or materials
used and the criteria adopted to assess the deliverables. Summarize the solution,
the findings and the conclusion supported by available evidence. Identify the
actual development stage and distinguish technical checks from evaluation with
people or observed effects in the intended setting. If the scope changed,
explain its consequence when needed to understand the outcome. Do not present
expected benefits as demonstrated findings or turn the abstract into a plan
for future activities. Write one coherent paragraph using objective language,
without bibliographic citations or information absent from the report. Replace
these instructions with an accurate account of your completed project before
submission. Check the keywords and ensure that both language versions describe
the same objective, procedures, findings and limitations, following degree
program requirements.}
\undfsetcontent{keywords}{applied project; final report; assessment of results}
